\subsection{Same Cardinality, Finite, Countable, and Uncountable}

\begin{definition} \label{def:cardinality-equivalence} \hlt{Same Cardinality}\\
For sets $A,B$, write $A\sim_{\#}B$ iff a bijection $f:A\to B$ exists. Then $A$ and $B$ have \hlt{the same cardinality}, or are \hlt{equivalent}.
\end{definition}

\begin{proposition} \label{prop:cardinality-equivalence-relation} \hlt{Same Cardinality Is an Equivalence Relation}\\
Let $A,B,C$ be sets.
\begin{enumerate}[label=(\roman*)]
\setlength{\itemsep}{0pt}
\item $A\sim_{\#}A$.
\item $A\sim_{\#}B\Rightarrow B\sim_{\#}A$.
\item $(A\sim_{\#}B\land B\sim_{\#}C)\Rightarrow A\sim_{\#}C$.
\end{enumerate}
Consequently, $\sim_{\#}$ is an equivalence relation among sets.
\begin{proof}
\begin{enumerate}[label=(\roman*)]
\setlength{\itemsep}{0pt}
\item The identity $\operatorname{id}_A:A\to A$ is bijective, so $A\sim_{\#}A$.
\item Suppose $A\sim_{\#}B$ and choose a bijection $f:A\to B$. The inverse $f^{-1}:B\to A$ exists, and $f$ is its two-sided inverse.\\
Hence $f^{-1}$ is bijective by Theorem \ref{thm:bijective-iff-invertible-function}(i), so $B\sim_{\#}A$.
\item Suppose $A\sim_{\#}B$ and $B\sim_{\#}C$, and choose bijections $f:A\to B$ and $g:B\to C$.\\
The composite $g\circ f:A\to C$ is bijective by Proposition \ref{prop:composition-restriction-calculus}(vi), so $A\sim_{\#}C$.
\end{enumerate}
\hlt{Structural Commentary}: Identity, inversion, and composition verify the three equivalence-relation laws.\\
\hlt{Key Insight}: Bijections are closed under reversal and composition.
\end{proof}
\end{proposition}

\begin{theorem} \label{thm:schroeder-bernstein} \hlt{Schroeder--Bernstein Theorem}\\
Let $A,B$ be sets, and let $f:A\to B$ and $g:B\to A$ be injections. Then $A\sim_{\#}B$.
\begin{proof}
Set $A_0\defeq A\setminus g(B)$ and recursively set $B_n\defeq f(A_n)$ and $A_{n+1}\defeq g(B_n)$ for $n\in\N_0$.\\
Put $U_A\defeq\bigcup\limits_{n\in\N_0}A_n$, $U_B\defeq\bigcup\limits_{n\in\N_0}B_n$, $A_*\defeq A\setminus U_A$, and $B_*\defeq B\setminus U_B$.\\
The restriction $f_*\defeq f|_{U_A}:U_A\to U_B$ is bijective because $f$ is injective and $U_B=f(U_A)$. For $b\in B_*$, $g(b)\notin A_0$.\\
Suppose $g(b)\in A_{n+1}=g(B_n)$ with $n\in\N_0$. Then injectivity of $g$ gives $b\in B_n\subseteq U_B$, a contradiction. Thus $g(b)\in A_*$.\\
Conversely, if $a\in A_*$, then $a\notin A_0$, so $a\in g(B)$. Choose $b\in B$ satisfying $a=g(b)$. Suppose $b\in B_n$ with $n\in\N_0$.\\
Then $a\in A_{n+1}\subseteq U_A$, a contradiction. Hence $b\in B_*$. Therefore $g_*\defeq g|_{B_*}:B_*\to A_*$ is bijective.\\
Define $h:A\to B$ by $h(x)\defeq f_*(x)$ for $x\in U_A$ and $h(x)\defeq g_*^{-1}(x)$ for $x\in A_*$. The domain parts $U_A,A_*$ and image parts $U_B,B_*$ are disjoint partitions, and both branches are bijective. Hence $h$ is bijective, so $A\sim_{\#}B$.\\
\hlt{Structural Commentary}: Forward layers are matched by $f$; the complementary remainders are matched by $g^{-1}$.\\
\hlt{Key Insight}: Alternate the injections until the unmatched remainder can be reversed.
\end{proof}
\end{theorem}

\begin{definition} \label{def:finite-countable-uncountable} \hlt{Finite, Countable, and Uncountable}\\
Let $A$ be a set.
\begin{enumerate}[label=(\roman*)]
\setlength{\itemsep}{0pt}
\item $A$ is \hlt{finite} iff $(\exists n\in\N_0)\,A\sim_{\#}J_n$.
\item $A$ is \hlt{infinite} iff it is not finite.
\item $A$ is \hlt{countable} iff $A\sim_{\#}\N$.
\item $A$ is \hlt{uncountable} iff it is neither finite nor countable.
\item $A$ is \hlt{at most countable} iff it is finite or countable.
\end{enumerate}
Countable sets are also called denumerable. Here \emph{enumerable} more generally describes the image of a sequence.
\end{definition}

\begin{proposition} \label{prop:natural-numbers-do-not-inject-into-finite-sets} \hlt{The Natural Numbers Are Infinite}
\begin{enumerate}[label=(\roman*)]
\setlength{\itemsep}{0pt}
\item $(\forall m\in\N_0)$, no injection $\N\to J_m$ exists.
\item The set $\N$ is infinite.
\item If $A$ is countable, then $A$ is infinite.
\item If $F$ is finite and $A\subseteq F$, then $A$ is not countable.
\end{enumerate}
\begin{proof}
\begin{enumerate}[label=(\roman*)]
\setlength{\itemsep}{0pt}
\item For $r\in\N$, let $Q_r$ assert that no injection $\N\to J_{r-1}$ exists. Since $J_0=\emptyset$, no map $\N\to J_0$ exists, so $Q_1$ holds.\\
Assume $Q_r$ and suppose $u:\N\to J_r$ is injective. If $r\notin u(\N)$, then $u:\N\to J_{r-1}$ contradicts $Q_r$.\\
Otherwise let $s\in\N$ be the unique index satisfying $u(s)=r$.\\
Define $\tau:\N\to\N\setminus\{s\}$ by $\tau(k)\defeq k$ for $k<s$ and $\tau(k)\defeq k+1$ for $k\geq s$.\\
The map $\tau$ is injective and avoids $s$, so $u\circ\tau:\N\to J_{r-1}$ is injective, again contradicting $Q_r$. Thus $Q_{r+1}$ holds.\\
Hence $(\forall r\in\N)\,Q_r$ by Theorem \ref{thm:mathematical-induction}, proving (i).
\item If $\N$ were finite, a bijection $J_m\to\N$ would have an inverse injection $\N\to J_m$, contradicting (i).
\item Let $A$ be countable. If $A$ were finite, composing a bijection $\N\to A$ with a bijection $A\to J_m$ would contradict (i).
\item If a countable $A$ were contained in a finite $F$, then $\N\to A\hookrightarrow F\to J_m$ would be injective, contradicting (i).
\end{enumerate}
\hlt{Structural Commentary}: Induction deletes the unique preimage of the final point of a finite target.\\
\hlt{Key Insight}: Removing one domain point leaves $\N$ injectively equivalent to itself.
\end{proof}
\end{proposition}

\begin{proposition} \label{prop:finite-sum-reindexing} \hlt{Finite Sums Are Invariant under Reindexing}\\
Let $n\in\N_0$, let $c:J_n\to\C$, and let $\sigma:J_n\to J_n$ be a bijection. Then $\sum\limits_{i=1}^{n}c(\sigma(i))=\sum\limits_{i=1}^{n}c(i)$.
\begin{proof}
For $r\in\N$, let $P_r$ be the assertion of the proposition with $n=r-1$. Then $P_1$ is the empty-sum identity.\\
Assume $P_r$, put $m\defeq r-1$, and fix $c:J_{m+1}\to\C$ and a permutation $\sigma$ of $J_{m+1}$.\\
Put $j\defeq\sigma^{-1}(m+1)$ and define $\theta:J_m\to J_{m+1}\setminus\{j\}$ by $\theta(i)\defeq i$ for $i<j$ and $\theta(i)\defeq i+1$ for $i\geq j$.\\
Then $\theta$ and $\sigma\circ\theta:J_m\to J_m$ are bijective. Since $\sigma(j)=m+1$, commutativity gives $\sum\limits_{i=1}^{m+1}c(\sigma(i))=c(m+1)+\sum\limits_{i=1}^{m}c(\sigma(\theta(i)))$.\\
Latter expression equals $c(m+1)+\sum\limits_{i=1}^{m}c(i)=\sum\limits_{i=1}^{m+1}c(i)$ by $P_r$. Thus $P_{r+1}$ holds. Hence $(\forall r\in\N)\,P_r$ by Theorem \ref{thm:mathematical-induction}, proving $P_{n+1}$.\\
\hlt{Structural Commentary}: Induction removes the position mapped to the final index and reindexes the remainder.\\
\hlt{Key Insight}: A permutation changes positions without changing the multiset of summands.
\end{proof}
\end{proposition}

\begin{definition} \label{def:finite-set-summation} \hlt{Summation over a Finite Set}\\
Let $X$ be finite, put $n\defeq\#X$, let $f:X\to\C$, and choose a bijection $g:J_n\to X$. Define $\sum\limits_{x\in X}f(x)\defeq\sum\limits_{i=1}^{n}f(g(i))$.\\
If $X=\emptyset$, then $n=0$ and this value is the empty sum $0$. If $k:J_n\to X$ is another bijection, put $\sigma\defeq g^{-1}\circ k$.\\
Then $\sigma$ is a permutation of $J_n$, $k=g\circ\sigma$, and the two listings give equal sums by Proposition \ref{prop:finite-sum-reindexing}.\\
\hlt{Key Insight}: Reindexing changes the listing, not the sum.
\end{definition}

\begin{proposition} \label{prop:finite-set-summation-laws} \hlt{Finite Set Summation Laws}\\
Let $X,Y$ be finite sets.
\begin{enumerate}[label=(\roman*)]
\setlength{\itemsep}{0pt}
\item \hlt{Well-definedness}: If $f:X\to\C$, then $\sum\limits_{x\in X}f(x)$ is independent of the bijection $J_{\#X}\to X$.
\item \hlt{Substitution}: If $h:Y\to X$ is a bijection and $f:X\to\C$, then $\sum\limits_{x\in X}f(x)=\sum\limits_{y\in Y}f(h(y))$.
\item \hlt{Linearity}: If $f,g:X\to\C$, then $\sum\limits_{x\in X}(f(x)+g(x))=\sum\limits_{x\in X}f(x)+\sum\limits_{x\in X}g(x)$. If $c\in\C$, then $\sum\limits_{x\in X}cf(x)=c\sum\limits_{x\in X}f(x)$.
\item \hlt{Disjoint unions}: If $X\cap Y=\emptyset$ and $f:X\cup Y\to\C$, then $X\cup Y$ is finite. The identity $\sum\limits_{z\in X\cup Y}f(z)=\sum\limits_{x\in X}f(x)+\sum\limits_{y\in Y}f(y)$ holds.
\item \hlt{Finite Fubini}: The set $X\times Y$ is finite. If $f:X\times Y\to\C$, then the following identities hold.\\
$\sum\limits_{x\in X}\sum\limits_{y\in Y}f(x,y)=\sum\limits_{(x,y)\in X\times Y}f(x,y)$, and $\sum\limits_{y\in Y}\sum\limits_{x\in X}f(x,y)=\sum\limits_{(x,y)\in X\times Y}f(x,y)$.
\end{enumerate}
\begin{proof}
\begin{enumerate}[label=(\roman*)]
\setlength{\itemsep}{0pt}
\item Let $g,k:J_{\#X}\to X$ be bijections and put $\sigma\defeq g^{-1}\circ k$. Then $k=g\circ\sigma$, so the two sums agree by Proposition \ref{prop:finite-sum-reindexing}.
\item Choose a bijection $g:J_{\#Y}\to Y$. The composite $h\circ g:J_{\#Y}\to X$ is bijective, so $\#Y=\#X$ by Proposition \ref{prop:finite-cardinality-well-defined}.\\
After expanding the finite-set sums, $\sum\limits_{x\in X}f(x)=\sum\limits_{i=1}^{\#Y}f(h(g(i)))=\sum\limits_{y\in Y}f(h(y))$ by part (i).
\item Put $m\defeq\#X$, choose a bijection $\gamma:J_m\to X$, and set $a_i\defeq f(\gamma(i))$ and $b_i\defeq g(\gamma(i))$ for $i\in J_m$, with $a_i=b_i=0$ otherwise.\\
After expanding the three finite-set sums, the first identity holds by Proposition \ref{prop:finite-sum-basic-laws}(ii).\\
After the analogous expansions, the scalar-multiple identity holds by Proposition \ref{prop:finite-sum-basic-laws}(iii).
\item Put $m\defeq\#X$ and $n\defeq\#Y$. If $m=0$ or $n=0$, then one set is empty and both conclusions follow directly.\\
Suppose $m,n\in\N$, and choose bijections $\alpha:J_m\to X$ and $\beta:J_n\to Y$.\\
Define $\delta:J_{m+n}\to X\cup Y$ by $\delta(i)\defeq\alpha(i)$ when $i\leq m$ and $\delta(i)\defeq\beta(i-m)$ when $i>m$.\\
The two branches have disjoint images and make $\delta$ bijective because $X\cap Y=\emptyset$. Thus $X\cup Y$ is finite.\\
The identity in (iv) holds by Proposition \ref{prop:finite-sum-basic-laws}(i) after expansion along $\delta$.
\item Put $m\defeq\#X$ and $n\defeq\#Y$. The set $X\times Y$ is finite with $\#(X\times Y)=mn$ by Theorem \ref{thm:multiplication-principle-finite-products}(iii).\\
If $m=0$ or $n=0$, all three sums equal $0$. Otherwise choose bijections $\alpha:J_m\to X$ and $\beta:J_n\to Y$.\\
Define bijections $\ell,r:J_{mn}\to X\times Y$ by $\ell((i-1)n+j)\defeq(\alpha(i),\beta(j))$ and $r((j-1)m+i)\defeq(\alpha(i),\beta(j))$ for $i\in J_m$ and $j\in J_n$.\\
The first identity in (v) holds by Proposition \ref{prop:finite-sum-basic-laws}(i) and part (i) after expansion along $\ell$.\\
The second identity in (v) holds by Proposition \ref{prop:finite-sum-basic-laws}(i) and part (i) after expansion along $r$.
\end{enumerate}
\hlt{Structural Commentary}: Reindexing, concatenation, and lexicographic listings transfer ordinary finite-sum laws to sets.\\
\hlt{Key Insight}: A finite-set sum depends on values, not listing order.
\end{proof}
\end{proposition}

\begin{definition} \label{def:sequence-as-enumeration} \hlt{Sequences as Enumerations}\\
Let $A$ be a set.
\begin{enumerate}[label=(\roman*)]
\setlength{\itemsep}{0pt}
\item A \hlt{sequence} in $A$ is a map $x:\N\to A$, written $(x_n)_{n\in\N}$ with $x_n\defeq x(n)$.
\item A sequence $x:\N\to A$ \hlt{enumerates} $A$ iff $x(\N)=A$.
\item A sequence $x:\N\to A$ is a \hlt{distinct enumeration} of $A$ iff $x:\N\to A$ is bijective.
\end{enumerate}
Thus $A$ is countable iff it has a distinct enumeration.
\end{definition}

\begin{proposition} \label{prop:integers-countable} \hlt{The Integers Are Countable}\\
The set $\Z$ is countable.
\begin{proof}
Define $f:\N\to\Z$ by $f(2k)\defeq k$ and $f(2k-1)\defeq 1-k$ for $k\in\N$.\\
Define $e:\Z\to\N$ by $e(z)\defeq 2z$ for $z\geq1$ and $e(z)\defeq 1-2z$ for $z\leq0$.\\
The even and odd branches partition $\N$, while the positive and nonpositive branches partition $\Z$.\\
Direct substitution gives $e\circ f=\operatorname{id}_{\N}$ and $f\circ e=\operatorname{id}_{\Z}$. Thus $f$ is bijective by Theorem \ref{thm:bijective-iff-invertible-function}(i), so $\Z$ is countable.\\
\hlt{Structural Commentary}: Parity separates the positive and nonpositive integer branches.\\
\hlt{Key Insight}: Interleave the two one-sided integer lists.
\end{proof}
\end{proposition}

\begin{theorem} \label{thm:infinite-subset-of-countable-is-countable} \hlt{Infinite Subsets of Countable Sets}\\
Let $A$ be countable and let $E\subseteq A$ be infinite. Then $E$ is countable.
\begin{proof}
Choose a bijection $x:\N\to A$, write $x_n\defeq x(n)$, and set $I_E\defeq\{n\in\N:x_n\in E\}$. The restriction $x|_{I_E}:I_E\to E$ is bijective. Hence $I_E$ is infinite, since otherwise this bijection would make $E$ finite. Each intersection $I_E\cap J_n$ is finite by Proposition \ref{prop:addition-subtraction-and-division-principles}(ii).\\
For $k\in\N$, set $M_k\defeq\{n\in I_E:\#(I_E\cap J_n)=k\}$. We prove $(\forall k\in\N)\,M_k\neq\emptyset$ by induction.\\
Let $m\defeq\min I_E$ by well-ordering. Then $I_E\cap J_m=\{m\}$, so $m\in M_1$. Suppose $M_k\neq\emptyset$ and choose $m\in M_k$.\\
The set $I_E\setminus J_m$ is nonempty; otherwise $I_E\subseteq J_m$ would be finite by Proposition \ref{prop:addition-subtraction-and-division-principles}(ii). Let $r\defeq\min(I_E\setminus J_m)$ by well-ordering.\\
Then $I_E\cap J_r=(I_E\cap J_m)\mathbin{\dot\cup}\{r\}$, so $\#(I_E\cap J_r)=k+1$ by Proposition \ref{prop:addition-subtraction-and-division-principles}(i). Thus $r\in M_{k+1}$; $(\forall k\in\N)\,M_k\neq\emptyset$ by Theorem \ref{thm:mathematical-induction}.\\
If $n<m$ belong to $I_E$, then $I_E\cap J_m$ is the disjoint union of $I_E\cap J_n$ and the nonempty set $I_E\cap(J_m\setminus J_n)$.\\
Hence $\#(I_E\cap J_n)<\#(I_E\cap J_m)$ by Proposition \ref{prop:addition-subtraction-and-division-principles}(i), so each $M_k$ is a singleton. Let $n_k$ be its unique element.\\
If $n\in I_E$ and $k\defeq\#(I_E\cap J_n)$, then $n\in M_k$, so $n=n_k$. Thus $k\mapsto n_k$ is a bijection $\N\to I_E$.\\
The composite $k\mapsto x_{n_k}$ is a bijection $\N\to E$, so $E$ is countable.\\
\hlt{Structural Commentary}: Finite initial-segment ranks canonically enumerate the indices that hit the subset.\\
\hlt{Key Insight}: Rank each selected index by the number of earlier selected indices.
\end{proof}
\end{theorem}

\begin{proposition} \label{prop:subsets-and-images-at-most-countable} \hlt{Subsets and Images of At Most Countable Sets}\\
Let $A$ be at most countable.
\begin{enumerate}[label=(\roman*)]
\setlength{\itemsep}{0pt}
\item If $B\subseteq A$, then $B$ is at most countable.
\item If $f:A\to Y$ is a function, then $f(A)$ is at most countable.
\end{enumerate}
\begin{proof}
\begin{enumerate}[label=(\roman*)]
\setlength{\itemsep}{0pt}
\item Suppose $B\subseteq A$. If $A$ is finite, then $B$ is finite by Proposition \ref{prop:addition-subtraction-and-division-principles}(ii).\\
If $A$ is countable, then $B$ is finite or, when infinite, countable by Theorem \ref{thm:infinite-subset-of-countable-is-countable}. Thus $B$ is at most countable.
\item Suppose $f:A\to Y$ and first assume $A$ is finite. Put $m\defeq\#A$, choose a bijection $e:J_m\to A$.\\
Set $I\defeq\{i\in J_m:(\forall j\in J_m)\,(j<i\Rightarrow f(e(j))\neq f(e(i)))\}$. The set $I$ is finite by Proposition \ref{prop:addition-subtraction-and-division-principles}(ii).\\
The map $i\mapsto f(e(i))$ on $I$ is injective by the definition of $I$. Fix $y\in f(A)$. The nonempty set $\{i\in J_m:f(e(i))=y\}$ has a least element by well-ordering, and that element belongs to $I$. Thus the map is also surjective onto $f(A)$, so $f(A)$ is finite.\\
Now suppose $A$ is countable. Choose a bijection $e:\N\to A$ and set $I\defeq\{i\in\N:(\forall j\in\N)\,(j<i\Rightarrow f(e(j))\neq f(e(i)))\}$.\\
The set $I$ is at most countable by (i). The map $i\mapsto f(e(i))$ on $I$ is injective by the definition of $I$.\\
Fix $y\in f(A)$. The nonempty set $\{i\in\N:f(e(i))=y\}$ has a least element by well-ordering, and that element belongs to $I$.\\
Thus the map is also surjective onto $f(A)$, so $f(A)$ is at most countable.
\end{enumerate}
\hlt{Structural Commentary}: Subsets restrict a listing; images retain the first index of each attained value.\\
\hlt{Key Insight}: Deleting entries or repetitions cannot increase cardinality.
\end{proof}
\end{proposition}

\begin{proposition} \label{prop:at-most-countable-iff-enumerable-image} \hlt{At Most Countable Sets Are Exactly Enumerable Images}\\
For a set $B$, the following are equivalent:
\begin{enumerate}[label=(\roman*)]
\setlength{\itemsep}{0pt}
\item $B$ is at most countable.
\item $B=\emptyset$ or a surjection $f:\N\to B$ exists.
\end{enumerate}
Thus a nonempty set is at most countable iff it is the image of a sequence.
\begin{proof}
\begin{enumerate}[label=(\roman*)]
\setlength{\itemsep}{0pt}
\item[$(i)$] $\Rightarrow(ii)$ If $B=\emptyset$, then (ii) holds. If $B$ is countable, any bijection $\N\to B$ is surjective.\\
If $B$ is finite and nonempty, put $m\defeq\#B$, choose a bijection $e:J_m\to B$.\\
Define $f:\N\to B$ by $f(n)\defeq e(n)$ for $n\leq m$ and $f(n)\defeq e(m)$ for $n>m$. Then $f$ is surjective.
\item[$(ii)$] $\Rightarrow(i)$ If $B=\emptyset$, then $B$ is finite. If $f:\N\to B$ is surjective, then $B=f(\N)$ is at most countable by Proposition \ref{prop:subsets-and-images-at-most-countable}(ii).
\end{enumerate}
\hlt{Structural Commentary}: A finite list can repeat its final term; a countable list already supplies a surjection.\\
\hlt{Key Insight}: At most countable sets are sequence images or the empty set.
\end{proof}
\end{proposition}

\begin{proposition} \label{prop:positive-integer-pairs-countable} \hlt{Pairs of Positive Integers Are Countable}\\
The set $\N\times\N$ is countable.
\begin{proof}
For $r\in\N_0$, put $t_r\defeq\frac{r(r+1)}{2}$. Define $p:\N\times\N\to\N$ by $p(i,j)\defeq t_{i+j-2}+i$.\\
Fix $(i,j)\in\N\times\N$ and put $r\defeq i+j-1$; then $t_{r-1}<p(i,j)\leq t_r$. The intervals $(t_{r-1},t_r]$ are disjoint for $r\in\N$.\\
Hence equality of two values of $p$ forces equality of the coordinate sums and then equality of the first, and therefore also the second, coordinates. Thus $p$ is injective. Fix $q\in\N$. The set $R_q\defeq\{r\in\N:q\leq t_r\}$ is nonempty because $t_q\geq q$.\\
Let $r\defeq\min R_q$ by well-ordering. Minimality gives $t_{r-1}<q\leq t_r$. Set $i\defeq q-t_{r-1}$ and $j\defeq r+1-i$.\\
Then $i,j\in\N$ and $p(i,j)=q$, so $p$ is surjective. Hence $p$ is a bijection and $\N\times\N$ is countable.\\
\hlt{Structural Commentary}: Triangular intervals separate the diagonals and recover both coordinates from each index.\\
\hlt{Key Insight}: Coordinate-sum diagonals occupy disjoint finite intervals.
\end{proof}
\end{proposition}

\begin{theorem} \label{thm:countable-union-of-countable-sets} \hlt{Countable Union of Countable Sets}\\
Let $(E_n)_{n\in\N}$ be countable sets and set $S\defeq\bigcup\limits_{n\in\N}E_n$. Then $S$ is countable.
\begin{proof}
For $n\in\N$, let $\mathcal{B}_n$ be the set of bijections $\N\to E_n$. Then $(\forall n\in\N)\,\mathcal{B}_n\neq\emptyset$.\\
Choose $(x_n)_{n\in\N}$ satisfying $(\forall n\in\N)\,x_n\in\mathcal{B}_n$ by the axiom of choice, and write $x_{n,k}\defeq x_n(k)$.\\
Choose a bijection $d:\N\to\N\times\N$ by Proposition \ref{prop:positive-integer-pairs-countable}. For $t\in\N$, write $d(t)=(n_t,k_t)$ and define $a:\N\to S$ by $a(t)\defeq x_{n_t,k_t}$.\\
Then $a(\N)=S$, so $S$ is at most countable by Proposition \ref{prop:at-most-countable-iff-enumerable-image}.\\
The countable set $E_1$ is contained in $S$, so $S$ is not finite by Proposition \ref{prop:natural-numbers-do-not-inject-into-finite-sets}(iv). Hence $S$ is countable.\\
\hlt{Structural Commentary}: Choice selects row enumerations, and a pairing bijection flattens their two indices.\\
\hlt{Key Insight}: A countable array becomes one sequence after pairing its indices.
\end{proof}
\end{theorem}

\begin{illustrationfigure}[H]
  \begin{tikzpicture}[
  >=Stealth,
  every node/.style={font=\small},
  entry/.style={draw=gray!55, fill=gray!5, rounded corners=1pt, minimum width=0.62cm, minimum height=0.40cm},
  active/.style={draw=orange!60!black, fill=orange!10},
  scan/.style={->, blue!70!black, line width=0.7pt},
  guide/.style={orange!75!black, line width=0.65pt},
  output/.style={draw=gray!55, fill=gray!5, rounded corners=1pt, inner xsep=5pt, inner ysep=4pt}
]
  \draw[guide] (1.70,-0.65) -- (0.85,-1.30);
  \draw[guide] (2.55,-0.65) -- (0.85,-1.95);
  \draw[guide] (3.40,-0.65) -- (0.85,-2.60);

  \foreach \row/\name in {1/E_1,2/E_2,3/E_3,4/E_4} {
    \node[anchor=east] at (0,-0.65*\row) {$\name$};
    \foreach \col in {1,...,5} {
      \node[entry] (x-\row-\col) at (0.85*\col,-0.65*\row) {$x_{\row\col}$};
    }
    \node[anchor=west, gray!65] at (4.55,-0.65*\row) {$\cdots$};
  }

  \foreach \row/\col in {1/1,1/2,2/1,1/3,2/2,3/1,1/4,2/3,3/2,4/1} {
    \node[entry, active] at (0.85*\col,-0.65*\row) {$x_{\row\col}$};
  }

  \node[output, anchor=west] (list) at (5.25,-1.45)
    {$x_{11},\,x_{12},\,x_{21},\,x_{13},\,x_{22},\,x_{31},\,x_{14},\,\cdots$};
  \draw[scan] (4.55,-1.45) -- (list.west);
\end{tikzpicture}

  \caption{Diagonal scan of countable rows into one surjective sequence.}
  \label{fig:countable-union-diagonal-scan}
\end{illustrationfigure}

Figure \ref{fig:countable-union-diagonal-scan} illustrates the scan used in Theorem \ref{thm:countable-union-of-countable-sets}.\\
Every indexed entry occurs on a finite diagonal; repeated values do not affect surjectivity.

\begin{corollary} \label{cor:at-most-countable-union-at-most-countable} \hlt{At Most Countable Unions}\\
Let $A$ be at most countable and let $(B_\alpha)_{\alpha\in A}$ be an indexed family of at most countable sets. Then $\bigcup\limits_{\alpha\in A}B_\alpha$ is at most countable.
\begin{proof}
Set $T\defeq\bigcup\limits_{\alpha\in A}B_\alpha$. If $T=\emptyset$, then the conclusion holds. Suppose $T\neq\emptyset$ and set $I\defeq\{\alpha\in A:B_\alpha\neq\emptyset\}$.\\
The set $I$ is nonempty and at most countable by Proposition \ref{prop:subsets-and-images-at-most-countable}(i).\\
Choose a surjection $e:\N\to I$ by Proposition \ref{prop:at-most-countable-iff-enumerable-image}, and put $\alpha_j\defeq e(j)$. For $j\in\N$, let $\mathcal{S}_j$ be the set of surjections $\N\to B_{\alpha_j}$.\\
Then $(\forall j\in\N)\,\mathcal{S}_j\neq\emptyset$ by Proposition \ref{prop:at-most-countable-iff-enumerable-image}. Choose $(b_j)_{j\in\N}$ satisfying $(\forall j\in\N)\,b_j\in\mathcal{S}_j$ by the axiom of choice, and write $b_{j,k}\defeq b_j(k)$.\\
Choose a bijection $d:\N\to\N\times\N$ by Proposition \ref{prop:positive-integer-pairs-countable}. For $t\in\N$, write $d(t)=(j_t,k_t)$ and define $a:\N\to T$ by $a(t)\defeq b_{j_t,k_t}$.\\
Then $a(\N)=T$, so $T$ is at most countable by Proposition \ref{prop:at-most-countable-iff-enumerable-image}.\\
\hlt{Structural Commentary}: Surjections index the nonempty pieces and their elements before a pairing map flattens both indices.\\
\hlt{Key Insight}: Empty pieces can be discarded; all remaining rows admit sequence images.
\end{proof}
\end{corollary}

\begin{corollary} \label{cor:countable-cartesian-products} \hlt{Countable Cartesian Products}\\
Let $A,B$ be sets.
\begin{enumerate}[label=(\roman*)]
\setlength{\itemsep}{0pt}
\item If $A$ and $B$ are at most countable, then $A\times B$ is at most countable.
\item If $A$ and $B$ are countable, then $A\times B$ is countable.
\end{enumerate}
\begin{proof}
\begin{enumerate}[label=(\roman*)]
\setlength{\itemsep}{0pt}
\item For $a\in A$, set $S_a\defeq\{a\}\times B$. The map $b\mapsto(a,b)$ is a bijection $B\to S_a$, so $(S_a)_{a\in A}$ is an indexed family of at-most-countable sets. Since $A\times B=\bigcup\limits_{a\in A}S_a$, the product is at most countable by Corollary \ref{cor:at-most-countable-union-at-most-countable}.
\item By (i), $A\times B$ is at most countable. Choose $b_0\in B$. The map $a\mapsto(a,b_0)$ is a bijection $A\to A\times\{b_0\}$, so $A\times B$ contains the countable subset $A\times\{b_0\}$. Thus $A\times B$ is not finite by Proposition \ref{prop:natural-numbers-do-not-inject-into-finite-sets}(iv), and it is therefore countable.
\end{enumerate}
\hlt{Structural Commentary}: Slicing proves at-most-countability; one fixed-coordinate slice rules out finiteness.\\
\hlt{Key Insight}: A product is a union of copies of either factor.
\end{proof}
\end{corollary}

\begin{illustrationfigure}[H]
  \begin{tikzpicture}[
  x=0.82cm,
  y=0.82cm,
  >=Stealth,
  every node/.style={font=\small},
  point/.style={circle, fill=black, inner sep=1.25pt},
  active/.style={circle, draw=orange!60!black, fill=orange!12, line width=0.55pt, inner sep=2.2pt},
  guide/.style={orange!75!black, line width=0.65pt},
  scan/.style={->, blue!70!black, line width=0.7pt},
  output/.style={draw=gray!55, fill=gray!5, rounded corners=1pt, inner xsep=5pt, inner ysep=4pt}
]
  \foreach \i in {1,...,5} {
    \node[anchor=east] at (0,-0.42*\i) {$\i$};
    \node[anchor=south] at (0.47*\i,-0.1) {$\i$};
  }
  \node[anchor=south] at (2.65,-0.1) {$\cdots$};

  \draw[guide] (0.94,-0.42) -- (0.47,-0.84);
  \draw[guide] (1.41,-0.42) -- (0.47,-1.26);
  \draw[guide] (1.88,-0.42) -- (0.47,-1.68);

  \foreach \row in {1,...,5} {
    \foreach \col in {1,...,5} {
      \node[point] at (0.47*\col,-0.42*\row) {};
    }
  }

  \foreach \row/\col in {1/1,1/2,2/1,1/3,2/2,3/1,1/4,2/3,3/2,4/1} {
    \node[active] at (0.47*\col,-0.42*\row) {};
  }

  \node[output, anchor=west] (list) at (3.10,-1.05)
    {$(1,1),\,(1,2),\,(2,1),\,(1,3),\,(2,2),\,(3,1),\,\cdots$};
  \draw[scan] (2.70,-1.05) -- (list.west);
\end{tikzpicture}

  \caption{Diagonal enumeration of $\N\times\N$ by finite $i+j$ diagonals.}
  \label{fig:countable-diagonal-enumeration}
\end{illustrationfigure}

Figure \ref{fig:countable-diagonal-enumeration} depicts the $\N\times\N$ pairing in Proposition \ref{prop:positive-integer-pairs-countable}.\\
Each pair lies on one finite coordinate-sum diagonal, which determines its interval of indices.

\begin{theorem} \label{thm:finite-tuples-over-countable-set-countable} \hlt{Finite Tuples over a Countable Set}\\
Let $A$ be countable. For $n\in\N$, set $T_n\defeq\{(a_1,\ldots,a_n):(\forall k\in J_n)\,a_k\in A\}$. Then $T_n$ is countable.
\begin{proof}
For $r\in\N$, let $P_r$ assert that $T_r$ is countable. The map $a\mapsto(a)$ is a bijection $A\to T_1$, so $P_1$ holds. Assume $P_r$.\\
The product $T_r\times A$ is countable by Corollary \ref{cor:countable-cartesian-products}(ii), and concatenation bijects $T_r\times A$ with $T_{r+1}$.\\
Thus $P_{r+1}$ holds, so $(\forall r\in\N)\,P_r$ by Theorem \ref{thm:mathematical-induction}.\\
\hlt{Structural Commentary}: Induction appends one coordinate and transfers countability across a binary product bijection.\\
\hlt{Key Insight}: Every fixed finite product is built one coordinate at a time.
\end{proof}
\end{theorem}

\begin{corollary} \label{cor:rationals-countable} \hlt{The Rationals Are Countable}\\
The set $\Q$ is countable.
\begin{proof}
Set $D\defeq\{(p,q)\in\Z\times\Z:q\neq0\}$ and define $\pi:D\to\Q$ by $\pi(p,q)\defeq\frac{p}{q}$.\\
The set $\Z\times\Z$ is countable by Proposition \ref{prop:integers-countable} and Corollary \ref{cor:countable-cartesian-products}(ii), so $D$ is at most countable by Proposition \ref{prop:subsets-and-images-at-most-countable}(i).\\
Since $\Q=\pi(D)$, the set $\Q$ is at most countable by Proposition \ref{prop:subsets-and-images-at-most-countable}(ii).\\
The inclusion $\Z\hookrightarrow\Q$ identifies a countable subset, so $\Q$ is not finite by Proposition \ref{prop:natural-numbers-do-not-inject-into-finite-sets}(iv), and therefore $\Q$ is countable.\\
\hlt{Structural Commentary}: Integer pairs form an at-most-countable parameter set whose quotient map covers $\Q$.\\
\hlt{Key Insight}: Fractions are images of nonzero-denominator integer pairs.
\end{proof}
\end{corollary}

\begin{samepage}
\begin{proofmethod} \label{meth:cardinality-proof-patterns} \hlt{Cardinality Proof Patterns}\\
Let $A,B$ be sets. Choose the route matching the target cardinality claim.
\begin{enumerate}[label=(\roman*)]
\setlength{\itemsep}{0pt}
\item \hlt{Same cardinality}: Construct $f:A\to B$, then verify that $f$ is injective and surjective.
\item \hlt{Countable}: Construct a bijection $\N\to A$. Alternatively, prove that $A$ is at most countable and not finite.
\item \hlt{Uncountable}: Prove $A\neq\emptyset$. Fix an arbitrary map $s:\N\to A$ and construct $a_s\in A\setminus s(\N)$.\\
Then no sequence surjects onto $A$, so $A$ is not at most countable by Proposition \ref{prop:at-most-countable-iff-enumerable-image}.
\end{enumerate}
\hlt{Key Idea}: Bijections certify equal cardinality; diagonal witnesses obstruct every proposed enumeration.
\end{proofmethod}
\end{samepage}
